\caption{CARBON architecture and prediction-based analyses. (\textbf{a}) Four sequence channels and aligned phyloP values form a 101-nt input. Profile logits are normalised by softmax and multiplied by $\exp(z)$, with $z$ from the count head. (\textbf{b}) Convolutions of widths 3, 5 and 7 produce 96 concatenated channels. SE weights rescale the unpooled features before fusion; an independent input projection supplies the residual. (\textbf{c}) A $96\rightarrow64$ projection feeds a bidirectional GRU with 64 units per direction. Its output is added to a $96\rightarrow128$ residual projection before layer normalisation and dropout. Three $1\times1$ convolutions ($128\rightarrow64\rightarrow32\rightarrow1$) form the profile head; global mean pooling and two linear maps ($128\rightarrow32\rightarrow1$) form the count head. (\textbf{d}) Profile-shape, squared-CDF and count terms jointly supervise prediction (Equation~\ref{eq:total_loss}). (\textbf{e}) IG-guided window selection and TF-MoDISco motif recovery lead to the Tomtom/CISBP-RNA comparison described in \supsect{S1.7}. (\textbf{f}) Sequence substitutions measure prediction sensitivity. Sequence fragments and curves illustrate the workflow; batch normalisation, hidden-layer activations and head dropout are specified in Methods and \supsect{S1.4}.}
